\chapter{Статичний аналіз та сигнатурне виявлення}
\chaptershort{Статичний аналіз та сигнатурне виявлення}
\label{chap:ch06}

Статичний аналіз --- дослідження зразка без його виконання --- є історично першим і досі базовим підходом до виявлення й вивчення злоякісного коду. Він дешевий, швидкий, безпечний (без запуску реального коду) і дозволяє досліджувати всі гілки програми, а не лише виконані. На статичному аналізі ґрунтується сигнатурне виявлення, що залишається першим рубежем більшості антивірусних засобів, а також поглиблене дослідження зразків аналітиками. Водночас статичний аналіз має принципові межі: нерозв'язність точного аналізу коду (розділ~\ref{chap:ch03}) і техніки ухилення --- пакування, обфускація, шифрування (розділ~\ref{chap:ch08}), --- що приховують справжній код. Розділ розглядає рівні статичного аналізу, сигнатурне виявлення й мову правил YARA, дизасемблювання й аналіз коду, виявлення за подібністю й кластеризацію сімейств, евристичний статичний аналіз, його межі, і завершується кейсом розроблення правила виявлення для сімейства ЗПК. Динамічний аналіз, що доповнює статичний, розглянуто в розділі~\ref{chap:ch07}.

\section{Місце й рівні статичного аналізу}\label{sec:ch06-levels}

Статичний аналіз утворює кілька рівнів зростаючої глибини (табл.~\thechapter.1), від швидкого поверхневого до трудомісткого аналізу коду. На практиці аналіз є ітеративним: поверхневі ознаки визначають пріоритет і напрям глибшого дослідження, а результати живлять правила виявлення (розділ~\ref{chap:ch11}) і реагування (розділ~\ref{chap:ch16}). Конвеєр статичного аналізу з ескалацією глибини наведено на рис.~\thechapter.1.

\begin{tikzfig}
\begin{tikzpicture}[font=\scriptsize,
  b/.style={draw,rounded corners=2pt,minimum height=6.5mm,text width=20mm,align=center},
  >={Stealth[length=1.5mm]}]
\node[b] (s) at (0,0) {Поверхневий (геші, рядки, імпорти)};
\node[b] (sig) at (3.3,0) {Сигнатури й подібність};
\node[b] (code) at (6.6,0) {Аналіз коду (CFG, семантика)};
\draw[->] (s) -- node[above,font=\tiny] {підозрілий} (sig);
\draw[->] (sig) -- node[above,font=\tiny] {пріоритетний} (code);
\draw[->] (s.south) to[out=-90,in=-90] node[below,font=\tiny] {чистий $\to$ пропуск} (sig.south);
\node[font=\tiny] at (3.3,-1.2) {зростання глибини й вартості $\to$};
\end{tikzpicture}
\figcap{Рис.~\thechapter.1 --- Конвеєр статичного аналізу з ескалацією глибини}
\end{tikzfig}

\tabhead{Таблиця~\thechapter.1 --- Рівні статичного аналізу злоякісного коду}

\begin{center}\footnotesize
\begin{tabularx}{\linewidth}{L{24mm} Y L{16mm}}
\toprule
Рівень & Що досліджує & Вартість \\
\midrule
Поверхневий & геші, рядки, метадані, імпорти, ентропія & дуже низька \\
Сигнатурний & збіг з відомими патернами & низька \\
Структурний & секції, ресурси, подібність, imphash & низька \\
Аналіз коду & дизасемблювання, граф потоку керування & висока \\
Семантичний & поведінкова модель, абстрактне виконання & дуже висока \\
\bottomrule
\end{tabularx}
\end{center}

\section{Поверхневий статичний аналіз}\label{sec:ch06-surface}

Поверхневий аналіз дає швидкі індикатори без розбору коду. \emph{Криптографічні геші} (SHA-256) однозначно ідентифікують відомий зразок і використовуються для пошуку в репутаційних базах; проте, як зазначено в розділі~\ref{chap:ch03}, зміна одного байта повністю змінює геш, тому вони непридатні для виявлення варіантів. \emph{Рядки} (текстові фрагменти: URL, шляхи, повідомлення, імена функцій) часто розкривають функціональність та індикатори, хоча в сучасному ЗПК зазвичай зашифровані. \emph{Метадані PE} (час компіляції, компілятор, підпис, версійна інформація) дають контекст і ознаки підробки. \emph{Імпорти й imphash} (розділ~\ref{chap:ch05}) підказують можливості й кластеризують сімейства. \emph{Ентропія секцій}~\eqref{eq:ch03-entropy} виявляє пакування. Поєднання цих ознак у ваговій оцінці дає перший фільтр, що відсіює очевидно підозрілі зразки для глибшого аналізу й підвищує ефективну базову частоту на наступних етапах (розділ~\ref{chap:ch03}).

Практична цінність поверхневого аналізу полягає не лише в самих ознаках, а в їх \emph{зіставленні} з контекстом. Наприклад, високе значення ентропії секції коду саме собою не є доказом злоякісності: багато легітимних інсталяторів і захищених комерційних програм так само спаковані або зашифровані. Вагомість ознаки визначає відхилення від очікуваного для цього класу файлів: PE-файл, що оголошує себе системною утилітою, але не має цифрового підпису, має час компіляції в майбутньому й імпортує лише \texttt{LoadLibrary} та \texttt{GetProcAddress} (ознака динамічного розв'язання імпортів для приховування можливостей), значно підозріліший за спакований, але належно підписаний інсталятор. Тому зрілі конвеєри поверхневого аналізу працюють не з абсолютними порогами, а з профілями нормальності й репутаційними даними, накопиченими з потоку мільйонів файлів.

Окремо відзначимо роль рядків як джерела \emph{слідчих} індикаторів. Навіть коли основне корисне навантаження зашифроване, у зразку часто залишаються непрямі сліди: фрагменти шляхів складання (що видають середовище розробника та іноді мову інтерфейсу), назви бібліотек часу виконання, рядки повідомлень про помилки, адреси командних серверів, шаблони мьютексів для запобігання повторному запуску. Ці артефакти живлять як сигнатури й правила YARA, так і атрибуцію кампаній (розділ~\ref{chap:ch12}). Саме тому автори ЗПК докладають зусиль до мінімізації рядків --- шифрують константи, збирають рядки на стеку побайтово, завантажують їх з ресурсів під час виконання, --- а аналітики, відповідно, розробляють методи їх відновлення (емуляція, деобфускація рядків).

\section{Сигнатурне виявлення}\label{sec:ch06-signatures}

\subsection{Принцип і типи сигнатур}

\term{Сигнатура} --- це формалізований опис ознак, наявність яких свідчить про належність зразка до відомої загрози. Сигнатурне виявлення зіставляє досліджуваний об'єкт із базою сигнатур. Типи сигнатур: \emph{байтові} (послідовності байтів, характерні для зразка); \emph{геш-сигнатури} (точні геші файлів чи їх фрагментів); \emph{структурні} (ознаки формату, секцій, імпортів); \emph{правила} (логічні комбінації умов, як у YARA). Класичні антивірусні рушії (наприклад, ClamAV~\cite{ch06:clamav}) підтримують кілька форматів сигнатур.

Порівняння типів сигнатур за покриттям і стійкістю наведено в табл.~\thechapter.2.

\tabhead{Таблиця~\thechapter.2 --- Типи сигнатур та їхні властивості}

\begin{center}\footnotesize
\begin{tabularx}{\linewidth}{L{20mm} Y L{18mm}}
\toprule
Тип & Принцип & Покриття варіантів \\
\midrule
Геш-сигнатура & точний геш файлу & жодного (точний збіг) \\
Байтова & послідовність байтів коду & низьке (ламається пакуванням) \\
Структурна & секції, імпорти, формат & помірне \\
Правило (YARA) & логіка над ознаками & помірне--високе \\
Семантична & модель поведінки з коду & високе \\
\bottomrule
\end{tabularx}
\end{center}

Принципова перевага сигнатур --- висока точність (низька частота хибних тривог) для \emph{відомих} загроз і висока швидкість. Принципова вада --- неспроможність виявляти нові й змінені зразки: як показано в розділі~\ref{chap:ch01}, за експоненційного зростання кількості зразків~\eqref{eq:ch01-growth} сигнатурне покриття принципово відстає. Крім того, поліморфізм і пакування (розділ~\ref{chap:ch08}) роблять байтові сигнатури неефективними. Тому сигнатури залишаються необхідним, але недостатнім першим рубежем, що доповнюється поведінковими й статистичними методами (розділи~\ref{chap:ch07}, \ref{chap:ch09}). Варто підкреслити, що сигнатурний підхід зберігає цінність саме завдяки високій точності: за величезних потоків файлів (розділ~\ref{chap:ch03}) дешевий і точний перший рубіж, що впевнено блокує відомі загрози, звільняє дорогі поведінкові й статистичні методи для аналізу решти. Тому раціональна архітектура захисту не протиставляє сигнатури й поведінку, а вибудовує їх каскадом із зростанням вартості й покриття.

\subsection{Мова правил YARA}

Де-факто стандартом гнучких сигнатур став YARA~\cite{yara,ch06:yara-doc} --- мова опису правил, що поєднує текстові й байтові патерни з логічними умовами. Правило YARA складається з розділів: метаданих (опис, автор, посилання), рядків (патерни --- текстові, шістнадцяткові з масками, регулярні вирази) й умови (логічний вираз над наявністю патернів, їх кількістю, позицією, а також властивостями файлу). Умова може враховувати формат PE, розмір, ентропію, зміщення патернів. Приклад логіки правила (схематично, без працездатного коду): «файл є PE \texttt{AND} містить щонайменше два з трьох характерних рядків \texttt{AND} має imphash зі списку \texttt{AND} ентропія секції коду $> 7$». Така комбінація різнорідних ознак дає баланс між покриттям (виявляє варіанти сімейства) і точністю (уникає хибних тривог).

Сила YARA --- саме в \emph{модульності} умови. Окреме правило може посилатися на модуль \texttt{pe} (щоб перевіряти характеристики заголовка, секцій, таблиці імпортів, цифрового підпису), модуль \texttt{math} (для обчислення ентропії довільного діапазону), модуль \texttt{hash} (для гешів фрагментів). Це дозволяє висловлювати ознаки різної природи в одному логічному виразі: наприклад, вимагати, щоб точка входу лежала в секції з аномально високою ентропією \emph{й} водночас таблиця імпортів містила характерний для сімейства набір функцій. Така виразність робить YARA мовою не лише антивірусних лабораторій, а й центрів реагування (розділ~\ref{chap:ch16}) та платформ аналізу загроз, де правила обмінюються між організаціями як машиночитані індикатори (розділ~\ref{chap:ch12}).

Водночас гнучкість YARA покладає відповідальність на автора правила. Поширена помилка початківця --- писати правило за \emph{одним} зразком, механічно копіюючи в нього довгі байтові послідовності, що насправді є фрагментами стандартної бібліотеки компілятора або ресурсами. Таке правило або не виявляє жодного варіанта, крім вихідного (надмірна специфічність), або спрацьовує на всіх програмах, зібраних тим самим компілятором (хибні тривоги). Зріле правило спирається на ознаки, що відображають \emph{задум} автора ЗПК (специфічну логіку, унікальні константи алгоритмів, характерні послідовності системних викликів), а не на артефакти інструментарію складання.

Якість правила вимірюють тими самими метриками, що й детектор (розділ~\ref{chap:ch03}): повнотою (частка зразків сімейства, що спрацьовують) і точністю (частка істинних серед спрацювань). Ризик надто специфічного правила --- пропуск варіантів (низька повнота); надто загального --- хибні тривоги на легітимних файлах (низька точність). Для оцінювання ймовірності випадкового збігу байтової сигнатури довжиною $n$ байтів на випадкових даних використовують оцінку

\begin{equation}
P_{\text{випадк}} \approx (L - n + 1)\cdot 256^{-n},
\label{eq:ch06-random}
\end{equation}

де $L$ --- розмір файлу. Формула~\eqref{eq:ch06-random} показує, чому короткі сигнатури (малі $n$) дають багато хибних збігів: кожен додатковий байт зменшує ймовірність випадкового спрацювання в 256~разів, тому надійні байтові патерни мають довжину щонайменше кількох байтів специфічного коду, а не загальних послідовностей.

\section{Дизасемблювання та аналіз коду}\label{sec:ch06-disasm}

Глибокий аналіз потребує відновлення логіки коду через \term{дизасемблювання} --- перетворення машинного коду на асемблерні інструкції, а далі --- на граф потоку керування (CFG) і, за можливості, псевдокод. Інструменти (IDA Pro~\cite{ch05:eagle}, Ghidra, Binary Ninja, радару-подібні) автоматизують цей процес, але він має принципові труднощі~\cite{ch06:andriesse}: розрізнення коду й даних (на архітектурі зі змінною довжиною інструкцій неможливе точно в загальному випадку); непрямі переходи (адреса обчислюється під час виконання); самомодифікований код; антидизасемблерні прийоми (розділ~\ref{chap:ch08}). Ці труднощі є проявом нерозв'язності (розділ~\ref{chap:ch03}). На практиці застосовують два базові алгоритми дизасемблювання: \emph{лінійний прохід} (послідовне декодування байтів від початку секції), простий, але легко збиваний вставленням «сміттєвих» байтів між інструкціями, і \emph{рекурсивний обхід} (декодування за потоком керування, від точки входу за переходами й викликами), точніший, але безсилий перед непрямими переходами. Сучасні інструменти поєднують обидва з евристиками розпізнавання функцій, тому результат дизасемблювання реального зразка завжди слід розглядати як \emph{обґрунтоване припущення}, а не беззаперечну істину --- особливо для навмисно обфускованого коду.

Центральним поданням для аналізу коду є \term{граф потоку керування} (CFG) --- орієнтований граф, вершини якого є базовими блоками (лінійними послідовностями інструкцій), а дуги --- можливими переходами (рис.~\thechapter.2). CFG дозволяє зрозуміти структуру програми, виявити цикли (наприклад, цикл шифрування у вимагачі), умовні гілки (перевірки середовища антианалізу), недосяжний код. На основі CFG будують граф викликів функцій, що відображає архітектуру програми й використовується для структурної подібності (BinDiff). Обфускація (розділ~\ref{chap:ch08}) цілеспрямовано ускладнює CFG (вставлянням зайвих гілок, непрозорими предикатами, перехопленням потоку керування), тому відновлення зрозумілого CFG з обфускованого коду є окремою задачею деобфускації.

\begin{tikzfig}
\begin{tikzpicture}[font=\scriptsize,
  n/.style={draw,rounded corners=2pt,minimum height=5mm,text width=20mm,align=center,font=\tiny,inner sep=2pt},
  >={Stealth[length=1.4mm]}]
\node[n] (a) at (0,0) {B0: перевірка середовища};
\node[n] (b) at (-2.1,-1.3) {B1: вихід (виявлено пісочницю)};
\node[n] (c) at (2.1,-1.3) {B2: розшифрування корисного навантаження};
\node[n] (d) at (2.1,-2.7) {B3: цикл запису файлів};
\node[n] (e) at (0,-4.0) {B4: завершення};
\draw[->] (a) -- (b) node[midway,above left,font=\tiny] {так};
\draw[->] (a) -- (c) node[midway,above right,font=\tiny] {ні};
\draw[->] (c) -- (d);
\draw[->] (d) to[out=-30,in=30,looseness=5] (d) node[right=5mm,font=\tiny] {цикл};
\draw[->] (d) -- (e);
\draw[->] (b) |- (e);
\end{tikzpicture}

\figcap{Рис.~\thechapter.2 --- Спрощений граф потоку керування фрагмента вимагача: умовне розгалуження за результатом антианалізної перевірки (B0) та цикл шифрування файлів (B3) як характерна структурна ознака}
\end{tikzfig}

Семантичний підхід до виявлення~\cite{ch06:christodorescu} будує абстрактну модель поведінки з коду, що стійкіша до поверхневих змін, ніж байтові сигнатури, бо відображає \emph{що} код робить, а не \emph{як} він записаний. Платформи бінарного аналізу~\cite{ch06:brumley} надають проміжні подання (intermediate representation) для автоматизованих перетворень і міркувань про семантику незалежно від конкретної архітектури процесора. Проте повний семантичний аналіз дорогий і обмежений обфускацією, тому застосовується вибірково --- для пріоритетних зразків і розроблення стійких правил. Поєднання статичного й динамічного підходів (гібридний аналіз) долає частину обмежень кожного~\cite{ch03:hybrid}.

\section{Виявлення за подібністю та кластеризація}\label{sec:ch06-similarity}

Оскільки ЗПК існує сімействами, ключовим є виявлення за \emph{подібністю} до відомих зразків, а не точним збігом. Міри подібності (розділ~\ref{chap:ch03}): нечіткі геші (ssdeep, TLSH~\cite{ch06:tlsh}), що дають близькі значення для схожих файлів; подібність множин ознак (імпортів, $n$-грамів~\cite{ch03:kolter}) за коефіцієнтом Жаккара~\eqref{eq:ch03-jaccard}; структурна подібність графів функцій (BinDiff~\cite{ch06:bindiff}). Для масштабованого порівняння мільйонів зразків застосовують хешування ознак (BitShred~\cite{ch06:jang}) і локально-чутливе хешування.

Ці міри різняться не лише обчислювальною вартістю, а й \emph{рівнем абстракції}, на якому фіксують подібність, і, відповідно, стійкістю до видів зміни зразка. Нечіткі геші працюють на рівні байтів і добре виявляють дрібні правки (виправлення, перекомпіляцію з незначними змінами), але легко обманюються перепакуванням або зміною секцій. Подібність множин імпортів абстрагується від розташування коду й стійкіша до перекомпіляції, проте нечутлива до внутрішньої логіки. Структурна подібність графів функцій працює на рівні семантики потоку керування й найстійкіша до поверхневих перетворень, але найдорожча й сама вразлива до спрямованої обфускації CFG (розділ~\ref{chap:ch08}). На практиці системи триажу комбінують кілька мір: дешеві (нечіткі геші, imphash) для грубого попереднього групування мільйонів зразків, дорогі (графова подібність) --- лише всередині попередньо сформованих кластерів. Такий багаторівневий підхід робить порівняння за подібністю масштабованим на промислових колекціях.

\term{Кластеризація} групує зразки за подібністю, виявляючи сімейства й кампанії. Це має подвійну цінність: автоматична атрибуція нового зразка до відомого сімейства (пришвидшує реагування) і виявлення нових сімейств (як щільних груп раніше невідомих зразків). Якість кластеризації оцінюють мірами чистоти кластерів відносно відомої розмітки. Для оцінювання відстані між зразками, поданими векторами ознак $\mathbf{x}, \mathbf{y}$, часто використовують косинусну подібність:

\begin{equation}
\mathrm{sim}(\mathbf{x},\mathbf{y}) = \frac{\mathbf{x}\cdot\mathbf{y}}{\|\mathbf{x}\|\,\|\mathbf{y}\|} \in [0;1],
\label{eq:ch06-cosine}
\end{equation}

що нечутлива до масштабу й добре працює для розріджених ознакових векторів (частоти $n$-грамів, наявність функцій). Кластеризація й подібність є містком між статичним аналізом і машинним навчанням (розділ~\ref{chap:ch09}), де ознаки подібності стають входом класифікаторів.

\section{Евристичний статичний аналіз та його межі}\label{sec:ch06-heuristics}

\term{Евристичний} статичний аналіз шукає не конкретні зразки, а \emph{ознаки підозрілості}, характерні для ЗПК узагалі: підозрілі комбінації імпортів, аномальні секції, високу ентропію, ознаки пакування, нетипову структуру. Евристики виявляють нові зразки, яких немає в базі сигнатур, але ціною вищої частоти хибних тривог (легітимні програми теж пакують, шифрують, використовують «підозрілі» функції), що знову приводить до проблеми базової частоти (розділ~\ref{chap:ch03}). Тому евристики калібрують і поєднують з іншими методами.

Межі статичного аналізу визначаються двома чинниками. \emph{Теоретичний}: нерозв'язність точного аналізу (розділ~\ref{chap:ch03}) --- будь-який статичний аналізатор дає наближення з неминучими помилками. \emph{Практичний}: техніки ухилення (розділ~\ref{chap:ch08}) --- пакування й шифрування приховують справжній код, роблячи статичний аналіз файлу на диску марним без розпакування. Саме тому статичний аналіз доповнюють динамічним (розділ~\ref{chap:ch07}): запуск у контрольованому середовищі розгортає справжній код у пам'яті, після чого його можна аналізувати статично (аналіз дампа пам'яті). Порівняння методів тестування безпеки застосунків (статичних і динамічних) систематизовано у вітчизняних дослідженнях~\cite{ch06:sast}, а застосування машинного навчання в статичному аналізі --- у~\cite{ch03:opirskyi-ml}.

\section{Кейс: розроблення правила виявлення для сімейства}\label{sec:ch06-case}

Розглянемо послідовність розроблення правила YARA для виявлення сімейства ЗПК (захисна задача: створення детектора на основі аналізу відомих зразків).

\emph{Крок 1. Збір і кластеризація зразків.} За подібністю~\eqref{eq:ch06-cosine} та imphash формуємо набір зразків одного сімейства й контрольний набір легітимних і сторонніх файлів.

\emph{Крок 2. Виділення спільних ознак.} Знаходимо ознаки, спільні для сімейства й відсутні в контрольному наборі: характерні рядки (навіть зашифровані константи), послідовності коду, imphash, структурні особливості. Пріоритет --- ознакам, що відображають семантику (вищий рівень «піраміди болю», розділ~\ref{chap:ch02}), бо вони стійкіші до варіацій.

\emph{Крок 3. Формулювання умови.} Будуємо логічну умову, що вимагає збігу кількох незалежних ознак (а не однієї), щоб забезпечити точність. За~\eqref{eq:ch06-random} вибираємо достатньо специфічні байтові патерни.

\emph{Крок 4. Оцінювання.} Вимірюємо повноту на наборі сімейства й точність на контрольному наборі. За потреби коригуємо баланс: послаблюємо умову для підвищення повноти або посилюємо для зниження хибних тривог. Критично перевірити правило на великому наборі легітимних файлів, щоб уникнути хибних спрацювань у продакшні.

\emph{Крок 5. Супровід.} Правило старіє в міру появи нових варіантів, тому потребує моніторингу ефективності й оновлення. Стійкіші правила, засновані на семантиці, живуть довше за ті, що спираються на поверхневі байти.

Кейс ілюструє ключовий принцип: ефективне сигнатурне виявлення --- це не механічний запис байтів одного зразка, а узагальнення спільних, семантично значущих ознак сімейства з балансуванням повноти й точності та постійним супроводом. Такий підхід поєднує переваги сигнатур (точність, швидкість) із частковим покриттям варіантів, але не замінює поведінкових методів для принципово нових загроз.

\section*{Висновки до розділу}
\addcontentsline{toc}{section}{Висновки до розділу}

Статичний аналіз --- дешевий, швидкий і безпечний базовий підхід до виявлення й вивчення ЗПК, що охоплює всі гілки коду й утворює кілька рівнів від поверхневого до семантичного. Поверхневий аналіз (геші, рядки, метадані, імпорти, ентропія) дає швидкі індикатори й перший фільтр. Сигнатурне виявлення точне й швидке для відомих загроз, але принципово відстає від нових і безсиле проти поліморфізму й пакування; мова YARA забезпечує гнучкі правила, що поєднують різнорідні ознаки, якість яких балансує повноту й точність~\eqref{eq:ch06-random}. Дизасемблювання й аналіз коду розкривають семантику зразка, але мають принципові труднощі, зумовлені нерозв'язністю, а семантичне виявлення стійкіше до поверхневих змін. Виявлення за подібністю й кластеризація~\eqref{eq:ch06-cosine} дозволяють працювати із сімействами, автоматично атрибутувати нові зразки й виявляти нові сімейства, слугуючи містком до машинного навчання. Евристичний аналіз виявляє невідомі зразки ціною хибних тривог. Межі статичного аналізу --- теоретична нерозв'язність і практичні техніки ухилення --- долаються поєднанням із динамічним аналізом. Головний висновок: статичний аналіз є необхідним, але недостатнім; його сила --- у поєднанні рівнів і узагальненні семантично значущих ознак, а не в механічному зіставленні байтів.

\selfcheck
\begin{enumerate}
\item Опишіть рівні статичного аналізу (табл.~\thechapter.1). Чому аналіз є ітеративним?
\item Які індикатори дає поверхневий аналіз? Чому криптографічні геші непридатні для виявлення варіантів?
\item Що таке сигнатура? Назвіть типи сигнатур і поясніть переваги й вади сигнатурного виявлення.
\item Опишіть структуру правила YARA. Чому поєднання різнорідних ознак дає баланс повноти й точності?
\item За формулою~\eqref{eq:ch06-random} поясніть, чому короткі байтові сигнатури дають багато хибних збігів. Обчисліть $P_{\text{випадк}}$ для $n=8$, $L=10^6$.
\item Які принципові труднощі має дизасемблювання? Як вони пов'язані з нерозв'язністю (розділ~\ref{chap:ch03})?
\item Що таке семантичне виявлення й чому воно стійкіше до поверхневих змін, ніж байтові сигнатури?
\item Які міри подібності застосовують для виявлення варіантів? Поясніть косинусну подібність~\eqref{eq:ch06-cosine}.
\item Яку подвійну цінність має кластеризація зразків? Як її якість оцінюють?
\item У чому перевага й вада евристичного статичного аналізу? Як проблема базової частоти впливає на його застосування?
\item Які теоретичні й практичні чинники обмежують статичний аналіз? Як їх долають поєднанням з динамічним?
\item Опишіть послідовність розроблення правила виявлення для сімейства (кейс підрозділу~\ref{sec:ch06-case}). Чому пріоритет надають семантично значущим ознакам?
\item Чому ефективне сигнатурне виявлення не є механічним записом байтів одного зразка?
\item Дослідницьке: розробіть правило виявлення для обраного сімейства ЗПК, оцініть його повноту й точність на незалежних наборах і дослідіть стійкість до змін зразків (розділ~\ref{chap:ch08}).
\item Дослідницьке: порівняйте нечіткі геші (ssdeep, TLSH) і структурну подібність (BinDiff) за точністю кластеризації сімейств і обчислювальною вартістю на великій колекції зразків.
\end{enumerate}
